% =====================================================================
%   附录 J：详细分析 (Detailed Analysis)
% =====================================================================

\section{详细分析}\label{app:detailed}

\subsection{执行结果的有效性分析}

\subsubsection{有效性分析的动机}

在人工智能与机器学习领域，模型的效能、精度与可靠性对于实际落地至关重要。
对多个模型进行评估，有助于理解它们各自的优势与局限，从而为不同任务挑选
最合适的模型提供更明智的决策支持。本节有效性分析旨在以系统化的方式，
区分不同模型在任务完成度、上下文长度约束、返回格式准确性、动作准确性
与任务限制等方面的表现差异。对这些性能维度的深入研究，不仅增进我们对
模型能力的认知，还为未来应用的优化与改进奠定基础。

\subsubsection{有效性分析的定义}

为开展全面的有效性分析，我们将结果划分为五个不同的类别：

\begin{itemize}
    \item \textbf{完成（Completed）}：指模型无论最终结果如何，均按指令
          成功完成了任务。
    \item \textbf{超出上下文长度（Context Limit Exceeded）}：指模型的长度
          受到 API 约束，主要出现在 \code{text-davinci} 模型上。
    \item \textbf{格式无效（Invalid Format）}：指模型虽收到清晰的指令，
          却未能按预期格式返回响应。
    \item \textbf{动作无效（Invalid Action）}：指模型返回的格式正确，但其
          动作超出了所允许的动作空间，或动作参数有误。
    \item \textbf{超出任务限制（Task Limit Exceeded）}：指任务达到了其
          终止条件（如超出规定的轮数）。
\end{itemize}

通过将结果划分到这些类别中，我们能够更清晰地了解每个模型的优势与面临的
挑战，从而有针对性地加以改进。

\subsubsection{模型的有效性分析}

在我们的评估中，对 27 个不同模型的有效性表现进行了细致的审视。除具有
API 上下文长度严格限制的 \code{text-davinci} 模型之外，其它模型的结果
主要落在"完成"、"格式无效"、"动作无效"和"超出任务限制"四类中。

\begin{figure}[t]
    \centering
    \includegraphics[width=0.95\linewidth]{validity_placeholder.png}
    \caption{
        \textbf{图 6：模型有效性分析}。"格式无效"、"动作无效"与"文本长度
        超出"是常见错误。"上下文长度超出"错误仅出现在 \code{text-davinci}
        模型上。
    }
    \label{fig:validity}
\end{figure}

由所展示的详细分析可见几大关键趋势。如图~\ref{fig:validity} 所示，图表
直观地呈现了不同模型在各类别上的有效率分布，从而支持我们得出若干富有
洞见的结论。

\subsection{若干发现}

\subsubsection{指令遵循能力至关重要}

根据表~\ref{tab:outcome-cat} 所呈现的数据，我们可对商用 API 类模型与开源
模型之间的性能差异得出若干重要观察。值得注意的是，在"格式无效"与
"动作无效"这两类错误上，开源模型面临的挑战相对更多。具体而言，开源模型
中"格式无效"占比 10.4\%，相较商用 API 类模型的 6.0\% 明显偏高；同样地，
"动作无效"占比 13.6\%，亦高于商用 API 类模型的 4.6\%。这些差异或许
反映了商用模型在鲁棒性与泛化能力上的优势，亦可能是模型设计与训练阶段
——尤其是指令遵循——的细致打磨所致。

\begin{table}[h]
    \centering
    \small
    \caption{
        \textbf{表 6：两类模型执行结果分布的对比}。
    }
    \label{tab:outcome-cat}
    \begin{tabular}{l c c c c c}
        \toprule
        \textbf{模型类别} & \textbf{完成} & \textbf{CLE} & \textbf{格式无效}
        & \textbf{动作无效} & \textbf{TLE} \\
        \midrule
        商用 API 类     & 61.5\% & 3.0\% & 6.0\%  & 4.6\%  & 24.9\% \\
        开源模型       & 39.1\% & 0.0\% & 10.4\% & 13.6\% & 36.9\% \\
        \bottomrule
    \end{tabular}
\end{table}

值得一提的是，即便是最佳的模型，有时也会忽视重要的指令。
尽管我们在数据库任务中已明确给出正确格式（参见正文~\ref{sec:intro}~节
中示例），即便是 \code{gpt-4} 仍偶尔无法正确响应。
例如，对于"将 \code{Team Information} 表中 \code{Princeton Tigers} 队伍的
容量修改为 \code{45,000}"的任务，模型在尝试解释思路与命令时，未能按规范
给出带 \code{Action} 标签的正确 SQL 语句。这可能源于模型在训练或对齐过程中
将某些输出模式"内化"，进而忽略了特定的任务指令。

\subsubsection{智能体规划的一致性与执行性}

智能体的一项基础能力是具备连贯统一的思维过程，从而基于现实条件制定并实施
可行的计划（参见图~\ref{fig:validity}）。
许多模型具备在遇到问题时进行分析并制定初始计划的能力。然而，即便是某些
最先进的模型，也容易偏离或遗忘其最初的计划。不同模型在执行计划时保持思维
序列一致性的能力差异较大。这种能力深刻影响着 LLM 作为智能体的效能与
运作潜力。此处以家务整理环境为例加以说明。

家务整理环境涵盖一个仿真的家庭场景，模型需基于任务所提供的环境观察与要完成的
目标，从给定的动作空间中选取合适的动作。环境中实体众多，可选动作丰富，
赋予了模型较高的自由度，同时也对其保持清晰连贯思维的能力提出了严峻的
挑战。

一个由 \code{gpt-4} 完成家务任务的示例（节选，原文见附录 J.2.2）：
\begin{quote}\small\ttfamily
\textbf{任务}：你正处在一个房间中央。环顾四周，你可以看到一个 4 号橱柜、
3 号橱柜、2 号橱柜、1 号橱柜、1 号台面、1 号垃圾桶、2 号毛巾架、1 号毛巾架、
2 号水槽、1 号水槽、1 号马桶、1 号卫生纸架、1 号毛巾架。
你的任务是：把一块干净的肥皂放到台面上。\\
\\
\textbf{成功示例}：\code{gpt-4} 通过检查库存→打开橱柜→找到肥皂→清洁→将肥皂
放到台面这一清晰的多步过程，顺利完成了任务。
\end{quote}

\subsubsection{任务完成的理解：完成轨迹中的 Token 与轮次分布}

本节我们将分析 LLM 完成任务时所需的 token 数与轮次分布，并对比不同模型的
差异。整体而言，开源模型在长任务上所需 token 数显著高于商用 API 模型，
而商用模型之间的差异则相对较小。这一现象进一步印证了开源模型在长任务
处理上的局限性。

\subsubsection{任务限制超出的主要原因：模型倾向于重复先前内容}

在对 \code{TLE}（任务限制超出）类失败案例的归因中，我们发现一个普遍现象——
模型倾向于在多轮交互中重复先前已生成的内容。这种"内容重复"并非源于问题本身
所需的迭代细化，而常常是因为模型陷入局部最优，难以跳出已生成的思路。
这一行为在多个开源模型中尤为显著，也部分解释了为什么开源模型在
家务整理与网页浏览等需要长期规划的任务上表现尤为乏力。

\subsubsection{代码微调对 LLM 作为智能体的影响}

从前文（4.3 节）所述的对比（\code{codellama} 与 \code{llama-2} 系列模型）
可以归纳出：代码微调的影响具有两面性。一方面，代码微调有助于模型在相对
程式化的任务（例如网页购物）上表现更佳；另一方面，它也可能削弱模型在
需要更广泛通用思考能力的任务上的表现，例如数字卡牌游戏与操作系统任务。
这一观察提示我们：在进行代码微调时，需要在"流程遵循能力"与"通用思考
能力"之间审慎权衡。

\subsubsection{自我纠正能力}

具备自我纠正能力是高级智能体的重要特征之一。然而，本文实验观察显示，
即便是 \code{gpt-4} 与 \code{claude-3} 等顶级模型，在自我纠正方面仍存在
显著局限。具体而言：（1）当模型已生成格式或语义错误的输出时，下一轮的
"自我纠正"往往只是对先前输出的字面级微调，而非结构性反思；（2）当环境
反馈明确指出前一步动作无效时，模型常常重复同一类型的错误动作。这表明，
当前 LLM 在"元认知"层面的能力仍不成熟，亟需进一步研究。
